\section{A composition obstruction to the standard projective tensor proposal}

\subsection{F-SM41: the truncated-disk rules do not define a bifunctor}

Work over a field \(k\) in the two-sort abelian category
\(\mathcal N\) of F-SM34, with
\(E(V)=V^{**}/V\).  Its projectives include
\[
 P_0(V)=(V,V,1,0),\qquad P_1(V)=(E(V),V,0,1).
\]
This finding concerns only the following proposed rules,
motivated by the tensor product of two-term cochain disks:
\[
 \begin{aligned}
 P_0(V)\star P_0(W)&=P_0(V\otimes W)\oplus P_1(V\otimes W),\\
 P_0(V)\star P_1(W)&=P_1(V\otimes W),\\
 P_1(V)\star P_0(W)&=P_1(V\otimes W),\\
 P_1(V)\star P_1(W)&=0.
 \end{aligned}                                                \tag{SM41.1}
\]
Use the standard disk splitting in which the two degree-one
coordinates are \((u,u+v)\).  On a cross map
\(a:P_0(V)\to P_1(W)\), defined by a linear map
\(V\to E(W)\), the proposed rule tensors by the forward
strength \(E(W)\otimes Z\to E(W\otimes Z)\).
On the cross map \(b:P_1(W)\to P_0(W)\) defined by
\(1_W\), use the standard first-factor disk rule.
These rules do not preserve composition, so they do not
define a bifunctor even on the indicated projective objects.

\begin{proof}
Choose \(W\) with \(E(W)\ne0\), for example
\(W=k^{(\mathbb N)}\), as proved in F-SM34 and F-SM36.
Fix \(0\ne e\in E(W)\).  Let
\[
 a:P_0(k)\longrightarrow P_1(W),\qquad
 b:P_1(W)\longrightarrow P_0(W)
\]
have component maps \(a=(1\mapsto e,0)\) and
\(b=(0,1_W)\), respectively.  They satisfy the two-sort
morphism relations, and their composite is zero in both
components:
\[
                              ba=0.                         \tag{SM41.2}
\]

Put \(T=P_0(k)\).  The normalized forward strength at
the scalar unit is the canonical identification
\(E(W)\otimes k=E(W\otimes k)\).  Hence the proposed
map on the first cross arrow is
\[
      a\star1_T=(a,0):P_0(k)\oplus P_1(k)\longrightarrow P_1(W).
                                                               \tag{SM41.3}
\]
For the second cross arrow, its degree-one coordinates
in the unsplit disk tensor are \((w,0)\).  In the chosen
coordinates \((u,u+v)\), these are \((u,v)=(w,-w)\).
Thus the standard rule is
\[
      b\star1_T=\binom{b}{-1}:
                  P_1(W)\longrightarrow P_0(W)\oplus P_1(W).
                                                               \tag{SM41.4}
\]
Multiplying these matrices gives
\[
 (b\star1_T)(a\star1_T)
       =\begin{pmatrix}0&0\\-a&0\end{pmatrix}\ne0.           \tag{SM41.5}
\]
The lower-left component is nonzero because \(a\ne0\);
this remains true in characteristic two.  A bifunctor
additive in the first variable would instead send
\((ba)\star1_T\) to zero by (SM41.2), contradicting
(SM41.5).  Neither an associativity nor a symmetry
coherence check is needed to exhibit the failure.
\end{proof}

\paragraph{Scope.}
The obstruction excludes these particular standard
truncated-disk maps together with the proposed object
rules and normalized forward strength.  It is not a
nonexistence theorem for every possible tensor product
on \(\mathcal N\).  In particular, it does not complete
the additive example into a counterexample to the full
handwritten symmetric-tensor hypotheses.
